% Options for packages loaded elsewhere
\PassOptionsToPackage{unicode}{hyperref}
\PassOptionsToPackage{hyphens}{url}
\PassOptionsToPackage{dvipsnames,svgnames*,x11names*}{xcolor}
%
\documentclass[
  10pt,
]{article}
\usepackage{lmodern}
\usepackage{amssymb,amsmath}
\usepackage{ifxetex,ifluatex}
\ifnum 0\ifxetex 1\fi\ifluatex 1\fi=0 % if pdftex
  \usepackage[T1]{fontenc}
  \usepackage[utf8]{inputenc}
  \usepackage{textcomp} % provide euro and other symbols
\else % if luatex or xetex
  \usepackage{unicode-math}
  \defaultfontfeatures{Scale=MatchLowercase}
  \defaultfontfeatures[\rmfamily]{Ligatures=TeX,Scale=1}
  \setmainfont[]{DejaVu Serif}
  \setmonofont[]{DejaVu Sans Mono}
\fi
% Use upquote if available, for straight quotes in verbatim environments
\IfFileExists{upquote.sty}{\usepackage{upquote}}{}
\IfFileExists{microtype.sty}{% use microtype if available
  \usepackage[]{microtype}
  \UseMicrotypeSet[protrusion]{basicmath} % disable protrusion for tt fonts
}{}
\makeatletter
\@ifundefined{KOMAClassName}{% if non-KOMA class
  \IfFileExists{parskip.sty}{%
    \usepackage{parskip}
  }{% else
    \setlength{\parindent}{0pt}
    \setlength{\parskip}{6pt plus 2pt minus 1pt}}
}{% if KOMA class
  \KOMAoptions{parskip=half}}
\makeatother
\usepackage{xcolor}
\IfFileExists{xurl.sty}{\usepackage{xurl}}{} % add URL line breaks if available
\IfFileExists{bookmark.sty}{\usepackage{bookmark}}{\usepackage{hyperref}}
\hypersetup{
  colorlinks=true,
  linkcolor=Maroon,
  filecolor=Maroon,
  citecolor=Blue,
  urlcolor=Blue,
  pdfcreator={LaTeX via pandoc}}
\urlstyle{same} % disable monospaced font for URLs
\usepackage[margin=0.6in,landscape]{geometry}
\usepackage{longtable,booktabs}
% Correct order of tables after \paragraph or \subparagraph
\usepackage{etoolbox}
\makeatletter
\patchcmd\longtable{\par}{\if@noskipsec\mbox{}\fi\par}{}{}
\makeatother
% Allow footnotes in longtable head/foot
\IfFileExists{footnotehyper.sty}{\usepackage{footnotehyper}}{\usepackage{footnote}}
\makesavenoteenv{longtable}
\setlength{\emergencystretch}{3em} % prevent overfull lines
\providecommand{\tightlist}{%
  \setlength{\itemsep}{0pt}\setlength{\parskip}{0pt}}
\setcounter{secnumdepth}{-\maxdimen} % remove section numbering
\renewcommand{\arraystretch}{1.6}

\author{}
\date{}

\begin{document}

\hypertarget{stage-1.1-liftoff-check-gravity-turn-loss-first-principles}{%
\section{Stage 1.1 --- Liftoff Check \& Gravity-Turn Loss (First
Principles)}\label{stage-1.1-liftoff-check-gravity-turn-loss-first-principles}}

\textbf{Era:} Redstone/Mercury-Redstone class, \textasciitilde1961
(fictional program), continuing Project Arrowhead. \textbf{Target time:}
\textasciitilde1.5 hours. \textbf{Prerequisite:} Stage 1.0 (Range Table
Targeting).

\hypertarget{why-this-problem-exists}{%
\subsection{Why this problem exists}\label{why-this-problem-exists}}

1.0 handed you Δv\_losses as a given number, the way a real FDO would
receive it from the trajectory design group. This problem \emph{is} that
group's work: where the number actually comes from, and --- as a bonus
you didn't ask for --- a check on whether the vehicle spec you were
handed even makes sense.

\hypertarget{part-1-sanity-check-the-vehicle}{%
\subsection{Part 1 --- Sanity-check the
vehicle}\label{part-1-sanity-check-the-vehicle}}

Before modeling anything about the ascent, check the most basic thing:
can this vehicle actually leave the pad? A rocket needs thrust greater
than its own weight at ignition, full stop.

\textbf{Candidate spec} (proposed by the vehicle design group): W₀ =
66,000 lbf, Wₚ = 40,000 lbf burned over t\_b = 143 s, Iₛₚ = 235 s.

\begin{verbatim}
ẇ = Wₚ / t_b                (weight flow rate)
T = ẇ · Iₛₚ                 (thrust, from the definition of Iₛₚ)
T/W₀ = liftoff thrust-to-weight ratio
\end{verbatim}

Compute T/W₀ for the original spec. A real booster of this class
typically lifts off around T/W₀ ≈ 1.2; anything at or below 1.0 cannot
lift off at all.

\textbf{What you should find:} T/W₀ ≈ 0.996. This vehicle cannot fly. It
was never caught because 1.0's calculation only used the
\emph{integrated} rocket equation (total Δv from the mass ratio), which
doesn't care about instant-by-instant thrust --- a completely reasonable
thing to miss if you never build a real time-history of the flight,
which is exactly what 1.0 didn't do and this problem does.

\hypertarget{part-2-corrected-data-and-the-real-ascent}{%
\subsection{Part 2 --- Corrected data, and the real
ascent}\label{part-2-corrected-data-and-the-real-ascent}}

The propulsion group reviewed the spec and traces the error to an
overstated liftoff weight. \textbf{Corrected value: W₀ = 55,000 lbf}
(Wₚ, t\_b, and Iₛₚ unchanged). Recompute T/W₀ and confirm it's now
reasonable (target: \textasciitilde1.2).

With the corrected vehicle, the guidance group ran a preliminary
trajectory-shaping study --- the kind of repetitive, closely-spaced
numerical integration that, even in 1961, nobody did by hand end-to-end.
That's the ``very limited computer'' moment for this problem: below is a
coarse extract from that run (7 points instead of the hundreds the real
integration would use), given to you the way a trajectory analyst would
actually receive it.

\textbf{Programmed pitch maneuver:} the guidance program holds the
vehicle vertical (γ = 90°) until t\_p = 13.0 s, then flies a prescribed
pitch profile that eases toward 15° by burnout. (The exact formula isn't
needed for this part --- it's given in the Verification step below, when
you build your own integrator.)

\textbf{Given --- trajectory checkpoints} (from the preliminary run):

\begin{longtable}[]{@{}llll@{}}
\toprule
t (s) & V (ft/s) & γ (°) & h (ft)\tabularnewline
\midrule
\endhead
0 & 0.0 & 90.0 & 0\tabularnewline
13 & 98.7 & 90.0 & 604\tabularnewline
39 & 499.8 & 46.53 & 6,485\tabularnewline
65 & 1,282.9 & 28.25 & 19,230\tabularnewline
91 & 2,372.1 & 20.57 & 37,878\tabularnewline
117 & 3,888.0 & 17.34 & 63,427\tabularnewline
143 & 6,284.2 & 15.98 & 100,235\tabularnewline
\bottomrule
\end{longtable}

\textbf{Given --- vehicle aerodynamic data} (from wind tunnel testing):
effective drag parameter Cd·A = 8.0 ft².

\textbf{Given --- atmosphere model:} isothermal exponential atmosphere
(the same simplified model Bate, Mueller \& White use for
atmospheric-drag perturbation work) ---

\begin{verbatim}
ρ(h) = ρ₀ · exp(−h/H)          ρ₀ = 0.0023769 slug/ft³, H = 23,800 ft
\end{verbatim}

\hypertarget{authentic-method-reference-material}{%
\subsection{Authentic method + reference
material}\label{authentic-method-reference-material}}

\textbf{Drag force:} D = ½·ρ(h)·V²·(Cd·A) --- standard aerodynamic drag
equation.

\textbf{Loss-rate bookkeeping.} At any instant, gravity is costing you
g₀·sin(γ) of your rocket's Δv-generating capability, and drag is costing
you (D/W)·g₀. This is genuine physics, not an artifact of this drill ---
it comes directly from the same equation of motion (dV/dt = c·ẇ/W −
D·g₀/W − g₀sinγ) that a full trajectory integration solves.

\textbf{The hand technique:} compute both loss rates at each of the
seven given checkpoints, then integrate (trapezoidal --- average the
rate at each pair of adjacent checkpoints, multiply by the time between
them, sum) to get total gravity loss and total drag loss over the burn.
This is \emph{not} the same as integrating the full coupled trajectory
yourself --- that's a much harder, numerically unstable problem by hand
(ask your instructor about the war story where the first attempt at this
diverged spectacularly with only five hand-sized steps). Using given
trajectory checkpoints and just integrating the loss \emph{rates}
sidesteps that entirely, and --- surprisingly --- still gets you a very
good answer, as you'll see in Verification.

\hypertarget{worksheet}{%
\subsection{Worksheet}\label{worksheet}}

\textbf{Part 1 --- original spec:}

\begin{longtable}[]{@{}ll@{}}
\toprule
Quantity & Value\tabularnewline
\midrule
\endhead
ẇ = Wₚ/t\_b &\tabularnewline
T = ẇ·Iₛₚ &\tabularnewline
T/W₀ (original, 66,000 lbf) &\tabularnewline
\bottomrule
\end{longtable}

\textbf{Part 2a --- corrected spec:}

\begin{longtable}[]{@{}ll@{}}
\toprule
Quantity & Value\tabularnewline
\midrule
\endhead
T/W₀ (corrected, 55,000 lbf) &\tabularnewline
\bottomrule
\end{longtable}

\textbf{Part 2b --- loss-rate table.} Fill in for each checkpoint.

\begin{longtable}[]{@{}lllllllll@{}}
\toprule
t (s) & V & h & γ & W(t) & ρ(h) & D & grav. rate = g₀sinγ & drag rate =
(D/W)g₀\tabularnewline
\midrule
\endhead
0 & 0.0 & 0 & 90.0° & & & & &\tabularnewline
13 & 98.7 & 604 & 90.0° & & & & &\tabularnewline
39 & 499.8 & 6,485 & 46.53° & & & & &\tabularnewline
65 & 1,282.9 & 19,230 & 28.25° & & & & &\tabularnewline
91 & 2,372.1 & 37,878 & 20.57° & & & & &\tabularnewline
117 & 3,888.0 & 63,427 & 17.34° & & & & &\tabularnewline
143 & 6,284.2 & 100,235 & 15.98° & & & & &\tabularnewline
\bottomrule
\end{longtable}

\textbf{Part 2c --- trapezoidal integration.}

\begin{longtable}[]{@{}llll@{}}
\toprule
Interval & Δt (s) & Avg grav. rate × Δt & Avg drag rate ×
Δt\tabularnewline
\midrule
\endhead
0--13 & 13 & &\tabularnewline
13--39 & 26 & &\tabularnewline
39--65 & 26 & &\tabularnewline
65--91 & 26 & &\tabularnewline
91--117 & 26 & &\tabularnewline
117--143 & 26 & &\tabularnewline
\bottomrule
\end{longtable}

Total gravity loss = \_\_\_\_\_\_\_\_\_\_ ft/s Total drag loss =
\_\_\_\_\_\_\_\_\_\_ ft/s Total Δv\_losses = \_\_\_\_\_\_\_\_\_\_ ft/s

Using Δv\_ideal = 9824 ft/s (from 1.0's Part A, same vehicle): V\_bo =
\_\_\_\_\_\_\_\_\_\_ ft/s

\end{document}
